\documentclass[12pt,a4paper]{article}
\usepackage[T2A,T1]{fontenc}
\usepackage[utf8]{inputenc}
\usepackage[english,russian]{babel}
\usepackage{XCharter}
\usepackage{amsmath,amssymb,mathtools}
\usepackage{array,longtable,booktabs}
\usepackage[margin=24mm]{geometry}
\usepackage[hidelinks]{hyperref}
\urlstyle{same}
\setlength{\parindent}{0pt}
\setlength{\parskip}{0.5em}
\setlength{\LTleft}{0pt}
\setlength{\LTright}{0pt}
\setlength{\tabcolsep}{4pt}
\emergencystretch=2em
\newcolumntype{P}[1]{>{\raggedright\arraybackslash}p{#1}}
\title{Как читать формулы Брусенцова: обозначения, операции и недосказанности}
\author{}
\date{}
\begin{document}
\maketitle

\emph{Разбор формул из публикаций 2001 и 2010 годов. Сопоставлены две статьи Брусенцова за 2010 год, глава 2005 года и статья 2003 года. Прямые пояснения автора отделены от межтекстового прочтения и нашей реконструкции.}

\section{Короткий ответ}

\textbf{В статье 2010 года не определены локально все символы и операции.} Перекрёстное чтение её с парной статьёй «Трехзначность отношения следования» (2010) помогает понять обозначения $I$ и $I'$, а глава 2005 года даёт подробное параллельное изложение статусов. Однако ни соседнее написание факторов, ни знак $\equiv$ не получают полного формального определения в статье с формулой. Эти пробелы нужно сохранять, а не заполнять молча.

Есть важное различие между двумя похожими записями. Формула
\[
V_x V_{x'} V_y V_{y'} \equiv 1
\]
описывает условие для первичных терминов и их противоположностей. Формула
\[
V_{xy} V'_{xy'} V_{x'y'}
\]
задаёт статусы трёх сочетаний признаков. Индексы показывают, что именно стоит после $V$: один термин или сочетание двух терминов.

\section{Формула из статьи 2010 года}

На первой странице PDF, соответствующей печатной странице 88, напечатано:
\[
(x\Rightarrow y) \equiv Axy \equiv I_{xy} I'_{xy'} I_{x'y'}.
\]
Формула сверена по изображению PDF и перекрёстно сопоставлена с диаграммой на третьей странице и другими работами. Последний множитель — именно $I_{x'y'}$: штрихи относятся к индексам $x'$ и $y'$, а оператор $I$ здесь без штриха. Это не запись $I'_{x'y'}$. В исходнике между $I$-выражениями нет явных знаков умножения или конъюнкции.

Статья вводит эту цепочку после слов о «необходимом следовании», а $Axy$ связывает с общеутвердительным суждением «Все $x$ суть $y$». На той же первой странице PDF автор также пишет:
\[
x\to y \equiv (x'\lor y) \equiv \neg(xy') \equiv I'_{xy'}.
\]
Внутри статьи это позволяет прочитать $I'_{xy'}$ как исключение сочетания $xy'$, контрпримера к утверждению «все $x$ суть $y$». Здесь $x\to y$ — материальная импликация в алгебре терминов Брусенцова, а не автоматически пропозициональная формула $P\to Q$ между двумя целыми высказываниями. Автор подчёркивает, что одного исключения $xy'$ для его необходимого следования недостаточно без принятого постулата сосуществования противоположностей.

В этой цепочке похожие знаки выполняют разные роли: $x'$ — противоположный термин, $\neg(xy')$ — отрицание сочетания $xy'$, а $I'_{xy'}$ — обозначение с отдельным штрихом у оператора $I$. Их нельзя сводить к одному универсальному знаку отрицания.

\textbf{Как разрешается расхождение в штрихах.} При низком разрешении маленький штрих у $x'$ в последнем индексе легко принять за штрих над $I$. Изображение первой страницы прямо показывает $I_{x'y'}$; диаграмма на PDF-с. 3 и параллельные профильные формулы в работах 2005, 2006 и 2011 годов подтверждают положительный статус сочетания $x'y'$ и привходящий статус $x'y$. Значит, здесь не обнаружена опечатка источника; обнаружена опасность типографического прочтения.

\subsection{Что формула сообщает, а чего не сообщает}

\begin{longtable}{P{0.17\textwidth}P{0.38\textwidth}P{0.35\textwidth}}
\toprule
\textbf{Обозначение} & \textbf{Что можно прочитать из статьи 2010 года} & \textbf{Статус пояснения} \\
\midrule
\endfirsthead
\toprule
\textbf{Обозначение} & \textbf{Что можно прочитать из статьи 2010 года} & \textbf{Статус пояснения} \\
\midrule
\endhead
$x$, $y$ & Термины/первичные признаки, связанные с классами вещей. & Дано в прозе. \\
$x'$, $y'$ & Противоположные признаки; штрих не объявлен в статье обычным логическим отрицанием высказывания. & Дано в прозе; знак $'$ не снабжён отдельным формальным определением. \\
$x\Rightarrow y$ & Отношение, которое автор называет необходимым следованием между терминами. & Название и словесный замысел даны; формальная семантика раскрывается через всю модель. \\
$x\to y$ & В формуле той же статьи — материальная импликация на уровне терминов/классовых выражений, преобразованная в $x'\lor y$ и $\neg(xy')$. & Не следует автоматически заменять пропозициональной формулой $P\to Q$ между целыми высказываниями. \\
$Axy$, $Ayx$ & Общеутвердительное суждение «Все $x$ суть $y$» и его обратное «Все $y$ суть $x$». & Понятно по контексту, но отдельная функция $A$ не задаётся словарной статьёй. \\
$I_{xy}$, $I'_{xy'}$, $I'_{x'y}$, $I_{x'y'}$ & Статья связывает штрихованные множители с исключением контрпримеров: $xy'$ для $Axy$ и $x'y$ для $Ayx$. Парная статья 2010 года поясняет $I_{xy}$ через возможность отношения «Некоторые $x$ суть $y$». & Это межтекстовое чтение, но не полное определение каждого множителя как простой формулы существования или отрицания. Сопоставление с $V$-записью — реконструкция, не объявленное автором тождество. \\
$V$, $V'$ & В статье 2001 года $V$ назван интегральной дизъюнкцией; штрих у оператора относится к антисуществованию/исключению сочетания. & Это специальные операции авторской модели; не отождествлять $V$ с обычным $\lor$, а $V'$ — с универсальным оператором отрицания. \\
$\lor$, $\neg$ & В формуле 2010 года входят в равенство $x\to y\equiv(x'\lor y)\equiv\neg(xy')$. & Здесь они применены к терминам и классовым выражениям; это не делает всю цепочку формулой высказываний с буквами $P,Q$. \\
$1$ & Правая часть условия $V_xV_{x'}V_yV_{y'}\equiv1$ в изложении 2001 года. & Контекстуально это единичное/утверждённое значение выражения в авторской алгебре; статья с формулой 2010 года его не использует и не задаёт для новичка. \\
Соседство множителей & В формуле факторы записаны подряд, например $I_{xy}I'_{xy'}$. & Для дальнейшего разбора можно принять совместное, конъюнктивное чтение; сама статья не формулирует это правило отдельно. \\
$\equiv$ & Соединяет выражения как тождественные или равносильные в рамках авторской записи; в $x\equiv y$ выражает тождество терминов. & Контекстуальное чтение; статья не задаёт формальную семантику знака для каждого случая. \\
\bottomrule
\end{longtable}

Таким образом, ответ на вопрос «определён ли каждый знак?» — \textbf{нет}. Перекрёстное чтение разрешает сомнение в последнем штрихе и даёт больше сведений о $I$, чем одна статья. Но формула всё ещё не полностью определена для нового читателя: статус каждого множителя, правило соседнего соединения и значение $\equiv$ требуют осторожного восстановления из контекста и других публикаций. Это недоопределённость изложения, а не бессмысленность формулы.

\subsection{Продолжение на второй странице статьи 2010 года}

На PDF-с. 2 (печатная с. 89) автор отмечает, что из \(x\Rightarrow y\) следует \(y'\Rightarrow x'\) (контрапозитивность), тогда как обратное \(y\Rightarrow x\) лишь возможно, но не необходимо. Затем он печатает:
\[
Axy\,I'_{x'y}
\equiv
AxyAyx
\equiv
I_{xy}I'_{xy'}I'_{x'y}I_{x'y'}
\equiv
(x\equiv y).
\]
Изображение страницы подтверждает штрихи: в условии \(I'_{x'y}\) штрих есть и у оператора \(I\), и у термина \(x'\); в последнем множителе произведения \(I_{x'y'}\) оператор не штрихован. Текст рядом поясняет, что для \(Ayx\) требуется исключить сочетание \(x'y\), а совместные \(Axy\) и \(Ayx\) выражают тождество \(x\equiv y\). Это проясняет, почему в формуле \(Axy\) на первой странице сочетание \(x'y\) не требуется исключать: оно исключается только при добавлении обратного отношения.

В соседнем тексте автор пишет \(Axy\Rightarrow Iyx\), описывая обратное отношение как возможное, но не необходимое. Это запись внутри алгебры терминов; статья не задаёт для этой стрелки отдельного правила в стиле пропозициональной импликации. Раскрытие длинной цепочки в произведение четырёх \(I\)-факторов сверено по изображению, но формальная операция соседнего соединения всё ещё не определена отдельным правилом.

\section{Более явная запись через $V$}

В статье «Трехзначная диалектическая логика» (2001, раздел о постулате сосуществования и универсуме Аристотеля) Брусенцов поясняет другую запись. Для первичных терминов $x$, $y$ и их противоположностей он записывает условие
\[
V_x V_{x'} V_y V_{y'} \equiv 1.
\]
Здесь $V$ автор называет символом «интегральной дизъюнкции»: это специальный оператор его модели, не автоматически обычная дизъюнкция $\lor$. Значения термина и его инверсии рассматриваются совместно. В статье это условие используется для постулата сосуществования противоположностей.

Для общеутвердительного отношения автор пишет:
\[
Axy \equiv V'_{xy'} V_x V_{x'} V_y V_{y'}
\equiv V_{xy} V'_{xy'} V_{x'y'}.
\]
Словесная расшифровка рядом с формулой сообщает: должны существовать вещи $xy$ и $x'y'$, а существование вещей $xy'$ исключено. Следующая фраза называет эту запись конъюнкцией существований и антисуществований. В этой записи:

\begin{itemize}
\item $V_x$, $V_{x'}$, $V_y$, $V_{y'}$ относятся к существованию первичных терминов и их противоположностей по отдельности;
\item $V_{xy}$ относится к существованию сочетания $xy$;
\item $V'_{xy'}$ отмечает исключённость/несуществование сочетания $xy'$;
\item $V_{x'y'}$ относится к существованию сочетания $x'y'$;
\item сочетание $x'y$ в этой формуле не указано, то есть его статус остаётся привходящим — оно может существовать или не существовать.
\end{itemize}

Здесь штрих на термине ($x'$) обозначает противоположный признак; штрих на операторе ($V'$) относится к антисуществованию/исключению сочетания. Это два разных места и два разных вклада в чтение формулы. В строке из статьи 2010 года эти позиции особенно легко спутать: $I'_{xy'}$ имеет штрих у оператора $I$, тогда как $I_{x'y'}$ имеет штрихи только у терминов индекса.

\textbf{Что даёт сопоставление двух статей 2010 года.} В короткой статье «Трехзначность отношения следования» символ $I_{xy}$ поясняется через возможность отношения «Некоторые $x$ суть $y$». Это помогает читать формулу длинной статьи, но не превращает её множители без остатка в обычные формулы существования. В статье 2003 года близкое обозначение $Ixy$ входит в описание классового суждения и соотносится с $Vxy$ при условии принятого универсума. Эти параллели подтверждают родство словаря, но не доказывают тождество семантики обозначений во всех публикациях.

Именно для $V$-формулы 2001 года автор прямо поясняет множители через существование и исключение сочетаний. Это локальное пояснение нельзя автоматически переносить на знак $+$ в профиле 2005 года или на все множители $I$ в формуле 2010 года. Там перевод статуса в утверждение «существует хотя бы один предмет» требует отдельного основания; для записи 2010 года это остаётся реконструкцией.

\subsection{Наглядное чтение на примере}

Пусть $x$ — «носит очки», а $y$ — «студент среди участников собрания». Если следовать словесному пояснению именно к $V$-формуле 2001 года, выражение $V_{xy} V'_{xy'} V_{x'y'}$ читается так:

\begin{itemize}
\item $V_{xy}$: в рассматриваемой области есть участники, которые носят очки и являются студентами;
\item $V'_{xy'}$: нет участников, которые носят очки, но не являются студентами;
\item $V_{x'y'}$: есть участники без очков, которые не являются студентами;
\item сочетание $x'y$ — участники без очков, но являющиеся студентами — этой формулой не запрещается и не требуется.
\end{itemize}

Это учебное пояснение комбинаций, а не пример из статьи Брусенцова. Оно передаёт прямую словесную расшифровку $V$-записи 2001 года. Для профиля $+-0+$ 2005 года и $I$-формулы 2010 года такое же экзистенциальное прочтение нельзя считать автоматическим: чтобы применять эти формулы к конечному списку, надо отдельно задать предметную область и обосновать перевод статусов в условия наличия или отсутствия предметов.

\section{Не смешивать условие универсума и профиль отношения}

Разница видна по индексам:

\begin{longtable}{P{0.39\textwidth}P{0.49\textwidth}}
\toprule
\textbf{Запись} & \textbf{Какой уровень она описывает} \\
\midrule
\endfirsthead
\toprule
\textbf{Запись} & \textbf{Какой уровень она описывает} \\
\midrule
\endhead
$V_xV_{x'}V_yV_{y'}\equiv1$ & Условие на область: каждый первичный термин и его противоположность имеет экземпляры. \\
$V_{xy}V'_{xy'}V_{x'y'}$ & Статусы трёх сочетаний в отношении $x,y$: наличие $xy$, исключение $xy'$, наличие $x'y'$. \\
$I_{xy}I'_{xy'}I_{x'y'}$ & Запись статьи 2010 года. Её параллель с $V$-формой правдоподобна, но сама статья не объявляет символы $I$ тождественными символам $V$. \\
\bottomrule
\end{longtable}

Условие универсума само по себе не говорит, какие именно совместные сочетания существуют. Переход к профилю ячеек делается внутри авторской модели, где приняты постулат сосуществования и соответствующие правила. Поэтому четыре одиночных условия $V_x,V_{x'},V_y,V_{y'}$ нельзя механически прочитать как утверждение о четырёх клетках $xy,xy',x'y,x'y'$.

\section{Практический порядок разбора формулы}

Когда встречается новая запись Брусенцова, читайте её в таком порядке:

\begin{enumerate}
\item Определите, что обозначают буквы: целые высказывания, первичные термины, предметы или сочетания признаков.
\item Посмотрите, что именно поставлено в индекс: один термин, например $V_x$, или клетка/сочетание, например $V_{xy}$.
\item Отделите штрих у термина от штриха у оператора. $x'$ — противоположный признак; $V'_{xy'}$ — исключённость клеточного сочетания в нотации источника.
\item Найдите пояснение операции соединения. Для $V$-формулы 2001 года автор говорит о конъюнкции; «интегральная дизъюнкция» $V$ остаётся специальным оператором. Для $I$-формулы 2010 года соседнее написание факторов похоже на совместное условие, но отдельное правило в статье не дано.
\item Установите условия модели: например, действует ли постулат сосуществования противоположностей.
\item Только после этого сравнивайте запись с материальной импликацией, обычной логикой классов или современной семантикой. Не переводите её в таблицу $P\to Q$ без отдельного обоснования.
\end{enumerate}

\section{Библиография}

\begin{itemize}
\item Н. П. Брусенцов. «Трехзначность отношения следования» (2010), печатные с. 86–87. \href{https://ternarycomp.cs.msu.ru/Papers/brusentsov_sborn2010_1.pdf}{PDF на сайте лаборатории троичной информатики МГУ}. \href{../sources/pdfs/Брусенцов_Трехзначность-отношения-следования_2010.pdf}{Локальная копия в корпусе}.
\item Н. П. Брусенцов. «Воссоздание аристотелевой безукоризненной силлогистики» (2010). \href{https://ternarycomp.cs.msu.ru/Papers/brusentsov_sborn2010_2.pdf}{PDF на сайте лаборатории троичной информатики МГУ}, электронные с. 1–4; печатные с. 88–91. \href{../sources/pdfs/Брусенцов_Воссоздание-аристотелевой-силлогистики_2010.pdf}{Локальная копия в корпусе}.
\item Н. П. Брусенцов. \href{https://www.computer-museum.ru/histussr/trilog11.htm}{«Трехзначная диалектическая логика»} (доложено в 2001 году; опубликовано в сборнике \emph{Программные системы и инструменты}, № 2, 2001, с. 36–44). Разделы о постулате сосуществования противоположностей и универсуме Аристотеля.
\item Н. П. Брусенцов. \href{https://computer-museum.ru/histussr/trilog0.htm}{«Трехзначная интерпретация силлогистики Аристотеля»} (2003). Архивная страница и производная PDF-копия в корпусе.
\end{itemize}

\end{document}
